\section{Primera Entrevista}
\begin{itemize}
  \item \underline{Alumnos:} En el sistema se desea controlar
  sus legajos y darles la oportunidad de generar formularios a
  los que tengan derecho. Es preferible colocarlos en un formato 
  de tabla, donde se puede ver cada uno de los componentes del
  legajo y si se halla entregado o no; a su vez que puede
  filtrarse por curso y por el estado de sus legajos (completo o
  incompleto). Todos los componentes podrían completarse de forma
  digital, excepto por el certificado médico; lamentablemente, en
  la actualidad sólo se tiene el modelo digital y se escriben los
  campos necesarios a mano. Lo mismo aplica para los formularios.
  Por último, es importante saber que se prefiere mantener la
  comunicación con los alumnos de forma personal; por lo cual no
  habrá un sistema de mensajería, además que ello lo otorga la
  plataforma Acadeu.
  \item \underline{Docentes:} El objetivo del sistema es gestionar
  sus legajos y habilitar formularios para que ellos completen.
  Se busca tener una representación de los legajos en una tabla
  que pueda filtrarse por si se completó el legajo o no. Los
  componentes de sus legajos son digitalizables, pero en la
  actualidad sólo el modelo es digital; ya que se siguen
  completando a mano y conservando el papel como único comprobante,
  tal como los formularios. Sin embargo, es destacable que el
  estado de los legajos se puede notificar a través del correo
  que provee el colegio.
  \item \underline{Exalumnos:} Existe un interés personal en conservar
  datos sobre los alumnos que se hayan recibido en el colegio.
  Principalmente, la información clave de ellos es sobre su egreso y
  formas de contacto.
\end{itemize}
\section{Segunda Entrevista}
\begin{itemize}
  \item \underline{Rectora:} Se encarga del manejo de asignación y
  relevo de docentes. Al principio de año, debe otorgarse cada uno
  de los cargos a los docentes registrados; no necesariamente será a
  todos, porque un docente puede no estar al principio de un año
  determinado. No existe ninguna consideración (más allá de los
  horarios) para asignar los puestos a los docentes, ya que la
  rectora tiene suficientes derechos para hacer las acciones que
  tiene asignadas. De ahí, en medio del año puede haber un relevo de
  docentes; lo cual puede significar una licencia de un docente, o
  una alta y baja del titular.
  \item \underline{Coordinación Escolar:} Se componen de
  coordinadores escolares y secretarios. Ellos se encargan de
  organizar los horarios, tanto al inicio del año como entremedio del
  mismo. El único horario que no tienen derechos para modificar dentro
  del año es Módulo. Por ahora, las asignaciones de horario y los
  cambios se hacen en una tabla de Excel; lo cual es muy inconveniente
  y refuerza la idea de implementar una verificación para los
  movimientos en el horario, considerando que es una operación de alta
  complejidad.
  \item \underline{Docente:} Según su salario, se clasifican en
  jornales y cargos; donde ambos roles pueden ser de un mismo docente.
  En cualquiera de los casos, debe cumplirse la cantidad máxima de 42
  horas del docente; incluyendo las horas trabajadas fuera del colegio.
  De ahí, cuando exista un cambio de horario, no se puede tener más de
  tres horas seguidas de una materia pedagógica (se excluye Módulo, el
  cual si tiene más de tres horas cátedra seguidas); además, el docente
  a cargo no debe tener horarios externos que le impidan ejercer el
  cargo.
  \item \underline{Alumno:} Cada uno recibe 3 calificaciones principales
  por materia de un trimestre, donde una materia que posea una duración
  menor a 3 meses debe seguir esta regla también (sólo para el trimestre
  en que transcurre); al finalizar el año, sólo podrán pasar aquellos que
  tengan a lo sumo 3 materias desaprobadas.
  \item \underline{Preceptor:} Se encarga de colocar la asistencia de
  docentes y alumnos en papel, donde el presente logra abarcar las horas
  contínuas que posea una materia; por cada uno de estos intervalos, es
  necesario realizar este registro.
\end{itemize}
